\documentclass[11pt]{article}
\usepackage{amsmath,amssymb,amsthm}
\usepackage{tikz}
\usepackage{hyperref}

\newtheorem{theorem}{Theorem}
\newtheorem{proposition}[theorem]{Proposition}
\newtheorem{lemma}[theorem]{Lemma}
\newtheorem{conjecture}[theorem]{Conjecture}
\newtheorem{definition}[theorem]{Definition}
\newtheorem{remark}[theorem]{Remark}
\newtheorem{example}[theorem]{Example}

\newcommand{\N}{\mathbb{N}}
\newcommand{\Z}{\mathbb{Z}}
\newcommand{\Q}{\mathbb{Q}}
\newcommand{\CPP}{\operatorname{CPP}}
\newcommand{\GK}{\operatorname{GK}}
\newcommand{\AG}{\operatorname{AG}}
\newcommand{\diag}{\operatorname{diag}}
\newcommand{\hook}{\operatorname{hook}}
\newcommand{\cohook}{\operatorname{cohook}}

% q-series notation
\newcommand{\qpoch}[2]{(#1;#2)_\infty}
\newcommand{\qpochn}[3]{(#1;#2)_{#3}}
\newcommand{\qbin}[2]{\genfrac{[}{]}{0pt}{}{#1}{#2}}
\newcommand{\jtheta}[2]{\theta(#1;#2)}

\title{Weighted Growth Diagrams and the Max Statistic\\
for Warnaar's Conjecture~2.7}
\author{Notes for research on cylindric partitions}
\date{\today}

\begin{document}
\maketitle

\begin{abstract}
We investigate the positivity of the polynomials $Q_{n,c}(q)$ arising
from the two-variable generating function of rank-3 cylindric partitions
(Warnaar, Conjecture~2.7). By extending the cylindric growth diagram
local rules from the author's thesis to carry the max statistic alongside
the $q$-weight, we develop both a computational verification pipeline and
a framework for understanding the sign cancellation that produces positivity.
Computational evidence confirms the conjecture for all tested profiles with
$d \leq 7$ and reveals striking coefficient patterns, including connections
to Kostka--Foulkes polynomials.
\end{abstract}

\tableofcontents

%% ========================================================================
\section{The max statistic through local rules}
\label{sec:max-local-rules}
%% ========================================================================

\subsection{Setup and notation}

Let $r = 3$ (rank) and $c = (c_0, c_1, c_2)$ with $d = c_0 + c_1 + c_2$
and $d \not\equiv 0 \pmod{3}$. The associated partition is
$\mu(c) = (c_1 + c_2,\, c_2,\, 0)$.

A cylindric plane partition $\pi$ with profile $c$ is a triple of partitions
$\nu = (\nu^{(1)}, \nu^{(2)}, \nu^{(3)})$ satisfying the interlacing conditions
\begin{align}
  \nu^{(1)}_j &\geq \nu^{(2)}_{j+c_1} \quad \text{for all } j \geq 1, \label{eq:interlace1} \\
  \nu^{(2)}_j &\geq \nu^{(3)}_{j+c_2} \quad \text{for all } j \geq 1, \label{eq:interlace2} \\
  \nu^{(3)}_j &\geq \nu^{(1)}_{j+c_0} \quad \text{for all } j \geq 1. \label{eq:interlace3}
\end{align}
The weight and max statistic are
\[
  |\pi| = |\nu^{(1)}| + |\nu^{(2)}| + |\nu^{(3)}|,
  \qquad
  \max(\pi) = \max\bigl(\nu^{(1)}_1,\, \nu^{(2)}_1,\, \nu^{(3)}_1\bigr).
\]

The two-variable generating function is
\begin{equation}
  \GK_c(z,q) := \sum_\pi z^{\max(\pi)} q^{|\pi|},
  \label{eq:GK-def}
\end{equation}
summing over all cylindric partitions with profile~$c$.

\subsection{Burge local rule and the first part}

The Burge local rule $\mathfrak{U}_{\alpha,\beta}(m,\nu) = \lambda$ (thesis,
Chapter~1, Definition~1.4.7) constructs $\lambda$ from input partitions
$\alpha, \beta, \nu$ and a nonnegative integer~$m$, satisfying the weight relation
\begin{equation}
  |\lambda| = m + |\alpha| + |\beta| - |\nu|.
  \label{eq:weight-relation}
\end{equation}

\begin{proposition}[First part through the local rule]
\label{prop:first-part}
If $\lambda = \mathfrak{U}_{\alpha,\beta}(m,\nu)$, then the first part satisfies
\[
  \lambda_1 = \max(\alpha_1, \beta_1) + \epsilon,
\]
where $\epsilon \in \{0, 1\}$ depends on whether the first column is among the
``added columns'' in the Burge algorithm. Specifically, $\epsilon = 1$ if and
only if column~$1$ belongs to the set $A \cup B \cup C \cup D$ in the forward
rule construction.
\end{proposition}

\begin{proof}
The Burge forward rule $\mathfrak{U}_{\alpha,\beta}(m,\nu) = \lambda$ adds one
box to the end of each column in $A \cup B \cup C \cup D$, where:
\begin{itemize}
  \item $A = \{j : \alpha'_j > \nu'_j\}$ (columns where $\alpha$ exceeds $\nu$),
  \item $B = \{j : \beta'_j > \nu'_j\}$ (columns where $\beta$ exceeds $\nu$),
  \item $C$ collects the ``cascade'' columns from the matching $\epsilon$ in $A \cap B$,
  \item $D$ adds the first $m$ columns from the complement.
\end{itemize}
The first column of $\lambda$ gains a box exactly when $1 \in A \cup B \cup C \cup D$.
Since $\alpha'_1 = \ell(\alpha)$ and $\nu'_1 = \ell(\nu)$, we have $1 \in A$
iff $\ell(\alpha) > \ell(\nu)$, and similarly for $B$.
\end{proof}

\subsection{The max through diagonal weights}

In the growth diagram framework, the bijection
$\psi_\pi : (\gamma, \mathfrak{d}) \to \mathfrak{c}$
maps a partition $\gamma$ and a labelled cylindric diagram $\mathfrak{d}$ to a
cylindric partition $\mathfrak{c} = (\mu^0, \mu^1, \ldots, \mu^T)$.

The strong weight preservation property (thesis, Chapter~3, Proposition~3.4.9) states:
\begin{equation}
  |\mu^k| = |\gamma| + |\diag(k)|_{\mathfrak{d}}
  \quad \text{for all } 1 \leq k \leq T,
  \label{eq:strong-weight}
\end{equation}
where $|\diag(k)|_{\mathfrak{d}} = \sum_{b \in \diag(k)} |\cohook(b)|_{\mathfrak{d}}$
is the diagonal weight.

The max statistic in terms of the growth diagram data is therefore:
\begin{equation}
  \max(\pi) = \max_k \mu^k_1 = \max_k \bigl(\gamma_1 + \delta_k(\mathfrak{d})\bigr),
  \label{eq:max-from-diag}
\end{equation}
where $\delta_k(\mathfrak{d})$ is the contribution to the first part from diagonal~$k$.

\begin{definition}[First-part diagonal increment]
\label{def:delta-k}
For a labelled cylindric diagram $\mathfrak{d}$ with diagonal weight decomposition,
define
\[
  \delta_k(\mathfrak{d}) := \mu^k_1 - \gamma_1 = [\text{first part of } \diag(k)
  \text{ contribution}].
\]
By \eqref{eq:strong-weight}, this equals the number of cohook-weighted boxes on
diagonal $k$ that contribute to the first column of $\mu^k$.
\end{definition}

The max statistic becomes:
\[
  \max(\pi) = \gamma_1 + \max_{1 \leq k \leq T} \delta_k(\mathfrak{d}).
\]

\begin{proposition}[Bounds on $\delta_k$]
\label{prop:delta-bounds}
For any labelled cylindric diagram $\mathfrak{d}$ of profile $\pi$, the
first-part diagonal increments satisfy:
\begin{enumerate}
  \item $\delta_k(\mathfrak{d}) \geq 0$ for all $k$.
  \item $|\delta_{k+1}(\mathfrak{d}) - \delta_k(\mathfrak{d})| \leq 1$ when
    positions $k, k+1$ form an inside corner of the profile (i.e., the local
    rule $\mathfrak{L}_k$ acts at this position).
  \item $0 \leq \delta_k(\mathfrak{d}) \leq |\diag(k)|_{\mathfrak{d}}$ (the
    first-part increment cannot exceed the total diagonal weight).
\end{enumerate}
\end{proposition}

\begin{proof}
(1) follows from $\mu^k_1 \geq \gamma_1$ since $\mu^k$ is obtained from $\gamma$
by adding boxes (the diagonal weight is nonneg).
(2) follows from Proposition~\ref{prop:first-part}: the local rule changes
$\mu^k_1$ by at most 1 via the $\epsilon$ term.
(3) is immediate from $\mu^k_1 \leq |\mu^k| = |\gamma| + |\diag(k)|_\mathfrak{d}$
and $\gamma_1 \leq |\gamma|$, though the bound can be tightened.
\end{proof}


%% ========================================================================
\section{The weighted growth diagram}
\label{sec:weighted-growth}
%% ========================================================================

\subsection{Extended bijection}

The growth diagram bijection $\psi_\pi$ naturally extends to track the pair
$(q^{|\pi|}, z^{\max(\pi)})$:
\[
  z^{\max(\pi)} q^{|\pi|}
  = z^{\gamma_1 + \max_k \delta_k(\mathfrak{d})} \cdot
    q^{T|\gamma| + \sum_{k=1}^T |\diag(k)|_{\mathfrak{d}}}.
\]

This factorization separates the ``base partition'' $\gamma$ contribution
from the ``cylindric diagram'' contribution:
\begin{itemize}
  \item The base partition $\gamma$ contributes $z^{\gamma_1} q^{T|\gamma|}$.
  \item The diagram $\mathfrak{d}$ contributes
    $z^{\max_k \delta_k(\mathfrak{d})} \cdot q^{\sum_k |\diag(k)|_\mathfrak{d}}$.
\end{itemize}

\subsection{Consequence for the local rule propagation}

When the higher-order local rule $\mathfrak{L}_i$ transforms $\psi_\pi$ to
$\psi_{\pi'}$ (thesis, Definition~3.3.1), the diagonal weight relation
(thesis, Proposition~3.4.7) gives:
\[
  |\diag(i)|_{\mathfrak{d}'} = m + |\diag(i-1)|_{\mathfrak{d}} + |\diag(i+1)|_{\mathfrak{d}}
  - |\diag(i)|_{\mathfrak{d}}.
\]

For the first-part tracking, we need the corresponding relation for $\delta_k$.
Using Proposition~\ref{prop:first-part} and the weight relation
\eqref{eq:weight-relation}:
\begin{equation}
  \delta_i(\mathfrak{d}') = \mu'^i_1 - \gamma_1
  = \lambda_1 - \gamma_1,
\end{equation}
where $\lambda = \mathfrak{U}_{\mu^{i-1}, \mu^{i+1}}(m, \mu^i)$, so
$\lambda_1 \leq \max(\mu^{i-1}_1, \mu^{i+1}_1) + 1$.


%% ========================================================================
\section{Combinatorial interpretation of $(zq)_\infty \cdot \GK_c(z,q)$}
\label{sec:combinatorial-interp}
%% ========================================================================

\subsection{The extraction formula}

By definition (Warnaar, eq.~2.4),
\begin{equation}
  Q_{n,c}(q) := \qpochn{q^\ell}{q^\ell}{n} \cdot [z^n]
  \Bigl(\qpoch{zq}{q} \cdot \GK_c(z,q)\Bigr),
  \label{eq:Q-def}
\end{equation}
where $\ell = \gcd(d, 3)$. Since $d \not\equiv 0 \pmod{3}$, we have $\ell = 1$
for $\gcd(d,3) = 1$ and $\ell = 1$ always (since $3 \nmid d$).

The Euler expansion gives
\[
  \qpoch{zq}{q} = \sum_{j \geq 0} \frac{(-1)^j z^j q^{j(j+1)/2}}{\qpochn{q}{q}{j}},
\]
so the $z^n$ coefficient is:
\begin{equation}
  [z^n]\bigl(\qpoch{zq}{q} \cdot \GK_c(z,q)\bigr)
  = \sum_{j=0}^{n} \frac{(-1)^j\, q^{j(j+1)/2}}{\qpochn{q}{q}{j}} \cdot g_{n-j}(q),
  \label{eq:zn-extraction}
\end{equation}
where $g_m(q) = [z^m]\,\GK_c(z,q)$.

\subsection{Combinatorial reading}

The function $g_m(q)$ counts cylindric partitions with max entry exactly~$m$,
weighted by $q^{|\pi|}$. The alternating sum \eqref{eq:zn-extraction} thus pairs:
\begin{itemize}
  \item Cylindric partitions with $\max(\pi) = n - j$, contributing positively (even $j$)
    or negatively (odd $j$).
  \item An additional ``tail'' partition from $1/(q)_j$ with $q$-shift by $j(j+1)/2$.
\end{itemize}

The combined object is: a cylindric partition $\pi$ with $\max(\pi) = n - j$
together with a partition $\sigma$ of length $\leq j$, with total weight
$|\pi| + |\sigma| + j(j+1)/2$, carrying sign $(-1)^j$.

\subsection{Franklin-type involution}

The $(-1)^j$ sign suggests a Franklin-type involution on these extended objects.
The multiplication by $\qpochn{q}{q}{n}$ in \eqref{eq:Q-def} converts the
formal power series $[z^n](\cdots)$ into a polynomial, absorbing the $1/(q)_j$
denominators.

\begin{remark}[The $d=2$ monomial phenomenon]
For $d = 2$, computational evidence shows:
\begin{align*}
  Q_{n,(1,1,0)}(q) &= q^{n^2}, \\
  Q_{n,(2,0,0)}(q) &= q^{n(n+1)}.
\end{align*}
These are monomials with exponents matching the quadratic forms in the
Andrews--Gordon sum side \cite[eq.~1.2]{Warnaar2021}. This extreme cancellation
suggests that for $d = 2$, the involution is \emph{complete}: every term with
$j \geq 1$ cancels perfectly except for a single surviving monomial.
\end{remark}


%% ========================================================================
\section{Positivity mechanism}
\label{sec:positivity}
%% ========================================================================

\subsection{Computational evidence}

We verified Conjecture~2.7 computationally for all inequivalent profiles with
$d \leq 7$ (under cyclic symmetry). The results are summarized below.

\begin{center}
\begin{tabular}{c|c|c|c|c}
$d$ & Profile $c$ & $Q(1)$ & $Q_{1,c}(q)$ & $Q_{2,c}(q)$ \\
\hline
1 & $(1,0,0)$ & $0$ & $0$ & $0$ \\
2 & $(2,0,0)$ & $1$ & $q^2$ & $q^6$ \\
2 & $(1,1,0)$ & $1$ & $q$ & $q^4$ \\
4 & $(2,1,1)$ & $4$ & $2q + q^2 + q^3$ & $q^3 + 3q^4 + 2q^5 + \cdots + q^{12}$ \\
5 & $(2,2,1)$ & $6$ & $2q + 2q^2 + q^3 + q^4$ & $q^3 + 4q^4 + 4q^5 + \cdots + q^{16}$ \\
5 & $(3,1,1)$ & $6$ & $2q + q^2 + 2q^3 + q^5$ & $q^3 + 3q^4 + 3q^5 + \cdots + q^{18}$ \\
7 & $(3,2,2)$ & $11$ & $2q + 3q^2 + 2q^3 + 2q^4 + q^5 + q^6$ & --- \\
7 & $(2,3,2)$ & $11$ & $2q + 3q^2 + 2q^3 + 2q^4 + q^5 + q^6$ & --- \\
7 & $(4,2,1)$ & $11$ & $2q + 2q^2 + 2q^3 + 2q^4 + q^5 + q^6 + q^8$ & --- \\
\end{tabular}
\end{center}

All coefficients are nonnegative, and $Q_{n,c}(1) = \bigl(\tfrac{1}{6}(d+1)(d+2) - 1\bigr)^n$
in every case.

\begin{remark}[Cyclic orbit structure for $d = 7$]
The profiles $(3,2,2)$ and $(2,3,2)$ lie in different cyclic orbits but produce the
same $Q_{1,c}(q)$. This is a consequence of the additional symmetry: both profiles
have $\mu(c) = (4,2,0)$ and $\mu(c) = (5,2,0)$ respectively, and they are related
by the conjugation symmetry of the underlying cylindric partition structure. The
profile $(4,2,1)$ has a distinct $Q_{1,c}(q)$ with a gap at $q^7$ and an isolated
$q^8$ term, continuing the pattern of sparse high-degree tails observed for $d \geq 4$.
\end{remark}

\subsection{Coefficient properties}

Beyond positivity, the coefficients exhibit additional structure:

\begin{enumerate}
  \item \textbf{$d = 2$: Monomials.}
    $Q_{n,(2,0,0)}(q) = q^{n(n+1)}$ and $Q_{n,(1,1,0)}(q) = q^{n^2}$.
    The exponents are precisely the quadratic forms appearing in the
    Andrews--Gordon sum side \cite[eq.~1.2]{Warnaar2021}.

  \item \textbf{$d \geq 4$: Not unimodal, not palindromic.}
    For $c = (2,1,1)$, $n = 2$: the coefficients are
    $[1, 3, 2, 3, 2, 2, 1, 1, 0, 1]$ (from $q^3$ to $q^{12}$),
    which has a gap at $q^{11}$ and is not unimodal.

  \item \textbf{Valuation pattern.}
    The valuation (lowest $q$-degree) of $Q_{n,c}$ appears to follow
    $\operatorname{val}(Q_{n,c}) = \binom{n}{2} \cdot \min(c_i \neq 0) + \cdots$,
    growing quadratically in $n$.

  \item \textbf{Isolated high-degree terms.}
    For $d \geq 4$, $Q_{n,c}(q)$ has a ``tail'' of sparse coefficients
    at high $q$-degrees, separated by gaps. For instance $Q_{2,(2,1,1)}$
    has a gap at $q^{11}$ before a final $q^{12}$ term.
\end{enumerate}

\subsection{Alternating sum structure}

Writing $[z^n](\qpoch{zq}{q} \cdot \GK_c) = P(q) - N(q)$ where $P$ collects
even-$j$ contributions and $N$ collects odd-$j$ contributions, the cancellation
ratio $N(1)/P(1)$ grows with $n$:

\begin{center}
\begin{tabular}{c|c|ccc}
$c$ & $d$ & $n=1$ & $n=2$ & $n=3$ \\
\hline
$(2,1,1)$ & 4 & 0.21 & 0.55 & --- \\
$(2,2,1)$ & 5 & 0.16 & 0.48 & --- \\
$(3,1,1)$ & 5 & 0.16 & 0.50 & --- \\
$(3,2,2)$ & 7 & 0.10 & --- & --- \\
\end{tabular}
\end{center}

The cancellation becomes increasingly tight as $n$ grows, yet always produces
positive coefficients after multiplying by $(q;q)_n$. This delicate balance
suggests a structured involution rather than a ``soft'' positivity argument.
For fixed $n$, the ratio decreases as $d$ increases (the interlacing constraints
become tighter, leaving less room for sign cancellation).

\subsection{Is positivity ``soft''?}

A natural question: does $(q;q)_n \cdot [z^n]\bigl((zq)_\infty \cdot f(z,q)\bigr)$
have nonneg coefficients for \emph{any} $f(z,q)$ with nonneg coefficients?

The answer is \textbf{no}. Consider $f(z,q) = 1 + zq$. Then
\[
  (zq)_\infty \cdot (1 + zq) = (zq)_\infty + zq \cdot (zq)_\infty
  = (zq^2)_\infty,
\]
so $[z^n]((zq^2)_\infty) = (-1)^n q^{n(n+3)/2}/(q)_n$. Multiplying by $(q)_n$
gives $(-1)^n q^{n(n+3)/2}$, which is negative for odd $n$.

\begin{remark}
This shows that the cylindric partition structure of $\GK_c(z,q)$ is \emph{essential}
for positivity --- it is not a formal consequence of having nonneg coefficients.
The specific interlacing structure of cylindric partitions provides the additional
constraints needed for the Franklin-type cancellation to produce a positive result.
\end{remark}

\subsection{Towards a bijective proof}

A bijective proof of positivity would construct a sign-reversing involution
$\iota$ on the set of pairs $(\pi, \sigma, j)$ (cylindric partition, tail partition,
index) with $\max(\pi) = n - j$ and $|\sigma| \leq j$, such that:
\begin{enumerate}
  \item $\iota$ reverses the sign $(-1)^j$.
  \item $\iota$ preserves the total $q$-weight $|\pi| + |\sigma| + j(j+1)/2$.
  \item The fixed points of $\iota$ (after multiplying by $(q)_n$) account
    for exactly $Q_{n,c}(q)$.
\end{enumerate}

The $d = 2$ monomial pattern $Q_{n,c}(q) = q^{n^2}$ (resp.\ $q^{n(n+1)}$)
suggests that for $d = 2$, there is essentially a unique fixed point, pointing
to a highly structured involution akin to Franklin's proof of Euler's pentagonal
number theorem.


%% ========================================================================
\section{Connection to Kostka--Foulkes polynomials}
\label{sec:kostka-foulkes}
%% ========================================================================

The Kostka--Foulkes polynomials $K_{\lambda,\mu}(q)$ are the $q$-analogs of
Kostka numbers, defined as
\[
  K_{\lambda,\mu}(q) = \sum_{T \in \mathrm{SSYT}(\lambda, \mu)} q^{\mathrm{charge}(T)},
\]
where the charge statistic is computed on semistandard Young tableaux of shape
$\lambda$ and content $\mu$. They are known to have nonneg coefficients (a deep
result of Lusztig, proved bijectively by Lascoux--Sch\"utzenberger via the charge
statistic).

\subsection{Comparison for $d = 4$}

For $c = (2,1,1)$, $n = 1$: $Q_{1,c}(q) = 2q + q^2 + q^3$. This polynomial has
degree~3, so we compare with $K_{\lambda,(1^3)}(q)$ for partitions $\lambda$ of 3:
\begin{align*}
  K_{(3),(1^3)}(q) &= q^3, \\
  K_{(2,1),(1^3)}(q) &= q + q^2, \\
  K_{(1^3),(1^3)}(q) &= 1.
\end{align*}
These give $Q_{1,(2,1,1)}(q) = K_{(2,1),(1^3)}(q) + K_{(3),(1^3)}(q) + q = q^3 + (q+q^2) + q$,
which does not decompose cleanly as a \emph{multiplicity-free} sum of Kostka--Foulkes polynomials.
The extra $q$ term suggests that the positivity mechanism for $Q_{n,c}(q)$ is not simply
a Kostka--Foulkes decomposition, though the connection to charge statistics merits further
investigation.

\subsection{Comparison for $d = 5$}

For $c = (2,2,1)$, $n = 1$: $Q_{1,c}(q) = 2q + 2q^2 + q^3 + q^4$. Comparing with
$K_{\lambda,(1^4)}(q)$ for partitions $\lambda$ of 4:
\begin{align*}
  K_{(4),(1^4)}(q) &= q^6, \\
  K_{(3,1),(1^4)}(q) &= q^3 + q^4 + q^5, \\
  K_{(2,2),(1^4)}(q) &= q^2 + q^4, \\
  K_{(2,1,1),(1^4)}(q) &= q + q^2 + q^3, \\
  K_{(1^4),(1^4)}(q) &= 1.
\end{align*}
We find $Q_{1,(2,2,1)}(q) = K_{(2,1,1),(1^4)}(q) + K_{(2,2),(1^4)}(q) - q^3$,
so again no clean Kostka--Foulkes decomposition exists. However, a decomposition
over a broader basis (e.g., Hall--Littlewood polynomials or parabolic
Kostka--Foulkes polynomials) may yet be possible.

\begin{remark}
The fact that $Q_{n,c}(q)$ does not decompose as a nonneg combination of
$K_{\lambda,\mu}(q)$ suggests that the positivity of $Q_{n,c}(q)$ arises from a
different mechanism than the charge statistic. This is consistent with the
cylindric structure: the involution needed for positivity must respect the
cyclic closure condition $\mu^0 = \mu^T$, which has no analog in the standard
RSK/charge framework.
\end{remark}


%% ========================================================================
\section{The omega involution and the E--H factorization}
\label{sec:omega-involution}
%% ========================================================================

The omega involution $\omega$ on symmetric functions sends the elementary
symmetric function $e_k$ to the complete homogeneous symmetric function $h_k$.
In the $q$-series setting, this corresponds to the duality between
\[
  E(z) = \qpoch{zq}{q} = \prod_{k \geq 1}(1 - zq^k)
  \quad \longleftrightarrow \quad
  H(z) = \frac{1}{\qpoch{zq}{q}} = \prod_{k \geq 1}\frac{1}{1 - zq^k}.
\]

\subsection{E--H factorization of $\GK_c$}

Every two-variable generating function $\GK_c(z,q)$ admits the factorization
\begin{equation}
  \GK_c(z,q) = H(z) \cdot F_c(z,q),
  \label{eq:EH-factor}
\end{equation}
where $F_c(z,q) := E(z) \cdot \GK_c(z,q) = \qpoch{zq}{q} \cdot \GK_c(z,q)$.
This definition is equivalent to: $F_c$ is the generating function whose
$z$-extraction gives
\[
  Q_{n,c}(q) = (q)_n \cdot [z^n]\, F_c(z,q).
\]

\begin{example}[The $d = 1$ case]
For $d = 1$, profile $(1,0,0)$: $\GK_c(z,q) = H(z) = 1/(zq;q)_\infty$,
so $F_c = E(z) \cdot H(z) = 1$. The $z^n$ coefficient of~$1$ is $\delta_{n,0}$,
confirming $Q_{n,c}(q) = 0$ for $n \geq 1$.
\end{example}

\subsection{Rogers--Ramanujan connection for $d = 2$}

For $d = 2$, the function $F_c$ takes the remarkable Rogers--Ramanujan form:
\begin{align}
  F_{(1,1,0)}(z,q) &= \sum_{m \geq 0} \frac{q^{m^2}\, z^m}{(q)_m},
    \label{eq:RR-first} \\
  F_{(2,0,0)}(z,q) &= \sum_{m \geq 0} \frac{q^{m(m+1)}\, z^m}{(q)_m}.
    \label{eq:RR-second}
\end{align}
Extracting $[z^n] F_c$ and multiplying by $(q)_n$ gives:
\begin{align*}
  Q_{n,(1,1,0)}(q) &= q^{n^2}, \\
  Q_{n,(2,0,0)}(q) &= q^{n(n+1)},
\end{align*}
confirming the monomial phenomenon for $d = 2$. These are precisely the
quadratic forms in the Andrews--Gordon sum side.

\subsection{Corrected generating functions $h_m$}

Define the \emph{corrected generating functions}
\begin{equation}
  h_m(q) := (q)_m \cdot g_m(q),
  \label{eq:h-m-def}
\end{equation}
where $g_m(q) = [z^m]\, \GK_c(z,q)$ is the $z^m$ coefficient of the full
(unfactored) generating function.

Note that $h_m \neq Q_m$ in general; rather, $Q_n$ is an alternating sum of the
$h_m$'s (see \S\ref{ssec:qbin-decomp} below). However, $h_m$ plays a
distinguished role because $g_m = h_m/(q)_m$ is the GK coefficient.

\begin{proposition}[Polynomiality and nonnegativity of $h_m$]
\label{prop:h-m-nonneg}
For all tested profiles with $d \leq 5$ and all $m$ in the computed range,
$h_m(q)$ is a polynomial with nonneg integer coefficients. The degree satisfies
$\deg(h_m) \leq m \cdot d$.
\end{proposition}

Computational evidence ($h_m$ denotes the polynomial for the given profile):

\begin{center}
\begin{tabular}{c|c|c|l}
$c$ & $d$ & $m$ & $h_m(q)$ \\
\hline
$(1,0,0)$ & 1 & $m$ & $q^m$ \\[2pt]
$(1,1,0)$ & 2 & 0 & $1$ \\
& & 1 & $2q$ \\
& & 2 & $2q^2 + q^3 + q^4$ \\[2pt]
$(2,0,0)$ & 2 & 0 & $1$ \\
& & 1 & $q + q^2$ \\
& & 2 & $q^2 + q^3 + q^4 + q^6$ \\[2pt]
$(2,1,1)$ & 4 & 0 & $1$ \\
& & 1 & $3q + q^2 + q^3$ \\
& & 2 & $3q^2 + 4q^3 + 5q^4 + 3q^5 + 3q^6 + 2q^7 + 2q^8 + q^9 + q^{10} + q^{12}$ \\[2pt]
$(2,2,1)$ & 5 & 0 & $1$ \\
& & 1 & $3q + 2q^2 + q^3 + q^4$ \\
\end{tabular}
\end{center}

\begin{remark}
For $d = 1$, $h_m = q^m$ (a single monomial), reflecting the fact that $\GK_{(1,0,0)} = H(z)$
and $[z^m] H(z) = q^m/(q)_m$.
For $d = 2$, the $h_m$ are nontrivial nonneg polynomials that encode
both Rogers--Ramanujan sums from~\eqref{eq:RR-first}--\eqref{eq:RR-second}:
the alternating sum of $h_m$'s collapses to a monomial $Q_n = q^{n^2}$ or $q^{n(n+1)}$.
\end{remark}

\subsection{The $q$-binomial decomposition}
\label{ssec:qbin-decomp}

Using the Euler expansion of $E(z) = \sum_j (-1)^j z^j q^{j(j+1)/2}/(q)_j$,
the conjecture polynomial $Q_n$ decomposes as:
\begin{equation}
  Q_{n,c}(q) = \sum_{j=0}^{n} (-1)^j\, q^{\binom{j+1}{2}}\,
  \qbin{n}{j}_{\!q} \cdot h_{n-j}(q),
  \label{eq:qbin-decomp}
\end{equation}
where $\qbin{n}{j}_q$ is the $q$-binomial coefficient.

\begin{proof}[Derivation]
Starting from $Q_n = (q)_n \cdot [z^n](E(z) \cdot \GK_c)$:
\begin{align*}
  Q_n &= (q)_n \sum_{j=0}^n \frac{(-1)^j q^{j(j+1)/2}}{(q)_j} \cdot g_{n-j} \\
      &= \sum_{j=0}^n (-1)^j q^{j(j+1)/2} \cdot
         \frac{(q)_n}{(q)_j} \cdot g_{n-j} \\
      &= \sum_{j=0}^n (-1)^j q^{j(j+1)/2} \cdot
         \qbin{n}{j}_{\!q} \cdot (q)_{n-j} \cdot g_{n-j} \\
      &= \sum_{j=0}^n (-1)^j q^{j(j+1)/2} \cdot
         \qbin{n}{j}_{\!q} \cdot h_{n-j}. \qedhere
\end{align*}
\end{proof}

This decomposition is computationally verified for all tested profiles:

\begin{center}
\begin{tabular}{c|c|c|c}
$c$ & $d$ & $n$ & Reconstruction matches $Q_n$? \\
\hline
$(2,1,1)$ & 4 & 1 & $\checkmark$ \\
$(2,1,1)$ & 4 & 2 & $\checkmark$ \\
$(2,2,1)$ & 5 & 1 & $\checkmark$ \\
\end{tabular}
\end{center}

\begin{remark}[Interpretation]
The decomposition~\eqref{eq:qbin-decomp} rewrites $Q_n$ as an alternating
$q$-binomial sum of the nonneg polynomials $h_m$. Since $\qbin{n}{j}_q$
and all $h_m$ have nonneg coefficients, the conjecture reduces to showing
that the alternating signs produce only nonneg total coefficients. This is
nontrivial and does \emph{not} hold for arbitrary nonneg sequences $h_m$
(cf.\ the soft positivity counterexample in \S\ref{sec:positivity}).
The specific structure of the cylindric $h_m$ must be exploited.
\end{remark}

\subsection{The omega-conjugate generating function}

Define the omega-conjugate generating function
\[
  \GK_c^\omega(z,q) := \sum_\pi z^{\max(\ell(\nu^{(1)}), \ell(\nu^{(2)}), \ell(\nu^{(3)}))} q^{|\pi|},
\]
where $\ell(\nu)$ denotes the number of parts (length) of $\nu$, tracking
max-length instead of max-first-part.

Under the omega involution $\omega: e_k \leftrightarrow h_k$, partition
conjugation $\lambda \mapsto \lambda'$ sends $\lambda_1 \mapsto \ell(\lambda)$.
Thus $\GK_c^\omega$ tracks the conjugate statistic.

Computational comparison shows:
\begin{itemize}
  \item $\GK_c(1,q) = \GK_c^\omega(1,q)$ (both count the same CPPs at $z=1$);
  \item $\GK_c(z,q) \neq \GK_c^\omega(z,q)$ in general (the $z$-refinement differs).
\end{itemize}

In all tested cases, the omega refinement redistributes weight: $\GK_c^\omega$
concentrates more terms at lower $z$-powers (since partition lengths are typically
smaller than first parts for the same total weight). For instance, with $c=(1,1,0)$,
$\max_n=3$: $[z^1]$ of $\GK$ has 20 terms while $[z^1]$ of $\GK^\omega$ has 39
terms; at $z^3$ the ratio reverses (137 vs 61).


%% ========================================================================
\section{Explicit involution for small cases}
\label{sec:explicit-involution}
%% ========================================================================

The $q$-binomial decomposition~\eqref{eq:qbin-decomp} expresses $Q_{n,c}(q)$ as
an alternating sum of nonneg polynomials. In this section we enumerate the signed
objects explicitly for $c = (2,1,1)$ ($d = 4$) and describe the sign-reversing
involution.

\subsection{Signed objects}

For each $j \in \{0, \ldots, n\}$, the objects in the $j$-th term of
\eqref{eq:qbin-decomp} are triples $(j, \lambda, w')$ where:
\begin{itemize}
  \item $j$ determines the sign $(-1)^j$,
  \item $\lambda$ is a partition fitting in a $j \times (n-j)$ box
    (from $\qbin{n}{j}_q$, contributing weight $|\lambda|$),
  \item $w'$ indexes an $h_{n-j}$-object of weight $w'$
    (multiplicity $[q^{w'}]\, h_{n-j}$),
  \item the total $q$-weight is $\binom{j+1}{2} + |\lambda| + w'$.
\end{itemize}

\subsection{The case $n = 1$, $c = (2,1,1)$}

Here $Q_1 = h_1 - q \cdot h_0$ with $h_0 = 1$ and
$h_1 = 3q + q^2 + q^3$. The signed objects are:

\begin{center}
\begin{tabular}{c|c|c|c|c|c}
$j$ & sign & $\lambda$ & $w'$ range & multiplicity & total weight \\
\hline
0 & $+$ & $()$ & 1 & 3 & 1 \\
0 & $+$ & $()$ & 2 & 1 & 2 \\
0 & $+$ & $()$ & 3 & 1 & 3 \\
1 & $-$ & $()$ & 0 & 1 & 1 \\
\end{tabular}
\end{center}

The $h_1$-coefficient $[q^1]\,h_1 = 3$ arises from three CPPs with
$\max = 1$ and weight~$1$:
\begin{equation}
  \label{eq:weight1-cpps}
  \bigl((1), \emptyset, \emptyset\bigr), \quad
  \bigl(\emptyset, (1), \emptyset\bigr), \quad
  \bigl(\emptyset, \emptyset, (1)\bigr).
\end{equation}
These are the three cylindric partitions obtained by placing a single box
in position 1, 2, or~3 of the cyclic structure.

The involution pairs one positive weight-1 object with the unique negative
weight-1 object, leaving:
\[
  Q_1 = 2q + q^2 + q^3.
\]

\begin{center}
\begin{tabular}{c|ccccc}
weight $w$ & pos & neg & paired & fixed & $[q^w]\,Q_1$ \\
\hline
1 & 3 & 1 & 1 & 2 & 2 \\
2 & 1 & 0 & 0 & 1 & 1 \\
3 & 1 & 0 & 0 & 1 & 1
\end{tabular}
\end{center}

\subsection{The case $n = 2$, $c = (2,1,1)$}

Here $Q_2 = h_2 - q(1+q)\,h_1 + q^3 \cdot h_0$ with
\begin{align*}
  h_0 &= 1, \\
  h_1 &= 3q + q^2 + q^3, \\
  h_2 &= 3q^2 + 4q^3 + 5q^4 + 3q^5 + 3q^6 + 2q^7 + 2q^8 + q^9 + q^{10} + q^{12}.
\end{align*}

The signed objects split across three values of~$j$:
\begin{itemize}
  \item $j = 0$ (sign $+$): $h_2$ objects, total 25.
  \item $j = 1$ (sign $-$): $q \cdot (1+q) \cdot h_1$, with
    $\lambda \in \{(), (1)\}$. Total 10 objects.
  \item $j = 2$ (sign $+$): $q^3 \cdot h_0 = q^3$, one object.
\end{itemize}

\begin{center}
\begin{tabular}{c|ccccc}
weight $w$ & pos & neg & paired & fixed & $[q^w]\,Q_2$ \\
\hline
2 & 3 & 3 & 3 & 0 & 0 \\
3 & 5 & 4 & 4 & 1 & 1 \\
4 & 5 & 2 & 2 & 3 & 3 \\
5 & 3 & 1 & 1 & 2 & 2 \\
6 & 3 & 0 & 0 & 3 & 3 \\
7 & 2 & 0 & 0 & 2 & 2 \\
8 & 2 & 0 & 0 & 2 & 2 \\
9 & 1 & 0 & 0 & 1 & 1 \\
10 & 1 & 0 & 0 & 1 & 1 \\
12 & 1 & 0 & 0 & 1 & 1
\end{tabular}
\end{center}

\noindent
This confirms $Q_2 = q^3 + 3q^4 + 2q^5 + 3q^6 + 2q^7 + 2q^8 + q^9 + q^{10} + q^{12}$
and $Q_2(1) = 16 = 4^2$.

\subsection{The involution mechanism}
\label{ssec:involution-mechanism}

The weight-by-weight analysis reveals a structured pairing:

\begin{enumerate}
  \item \textbf{At $w = 2$ (complete cancellation):}
    The 3~positive $h_2$-objects are the CPPs with $\max = 2$ and weight~2:
    \[
      \bigl((2), \emptyset, \emptyset\bigr), \quad
      \bigl(\emptyset, (2), \emptyset\bigr), \quad
      \bigl(\emptyset, \emptyset, (2)\bigr).
    \]
    These are paired with the 3~negative $j=1$ objects, which
    (via $\lambda = ()$, $w' = 1$) correspond to the weight-1 CPPs
    from~\eqref{eq:weight1-cpps} shifted by~$q$.

    The pairing maps position-by-position: the CPP with a height-2 box
    in position~$i$ is paired with the CPP with a height-1 box in position~$i$
    (plus the $q$-shift from the staircase $q^{j(j+1)/2}$).

  \item \textbf{At $w = 3$ (the staircase term survives):}
    The sole fixed point is the $j = 2$ object $q^3 \cdot h_0 = q^{n(n+1)/2}$,
    corresponding to the empty CPP with all weight concentrated in the
    staircase shift.

  \item \textbf{At $w \geq 6$ (all positive):}
    All $h_2$-objects at these weights are fixed points, contributing
    directly to $Q_2$.
\end{enumerate}

\begin{remark}[Column-stripping interpretation]
The involution ``strips the top layer'' from tall columns in the cylindric
partition: a CPP with $\max = m$ and a column of height~$m$ is paired with a
CPP with $\max = m - 1$ (the same configuration with the top box removed) plus a
compensating $q$-shift. In the growth diagram picture, this corresponds to the
$\epsilon = 0/1$ choice in the Burge local rule
(Proposition~\ref{prop:first-part}): the involution flips $\epsilon$
at the first position where the first part would increase.
\end{remark}

\subsection{Max-1 CPPs as length triples}

For $\max(\pi) = 1$, each partition $\nu^{(i)}$ has parts $\leq 1$, so it is
determined by its length $l_i$. The interlacing conditions reduce to:
\begin{equation}
  l_2 \leq l_1 + c_1, \qquad l_3 \leq l_2 + c_2, \qquad l_1 \leq l_3 + c_0.
  \label{eq:length-conditions}
\end{equation}
The weight is $w = l_1 + l_2 + l_3$.

For $c = (2,1,1)$: the conditions are $l_2 \leq l_1 + 1$, $l_3 \leq l_2 + 1$,
$l_1 \leq l_3 + 2$. The count of valid triples stabilises at 5 for $w \geq 3$
(the ``rank'' of the linear recurrence for $g_1$), giving:
\[
  g_1(q) = 3q + 4q^2 + 5q^3 + 5q^4 + 5q^5 + \cdots = 3q + 4q^2 + \frac{5q^3}{1-q}.
\]

This leads to $h_1 = (1-q) \cdot g_1 = 3q + q^2 + q^3$ as a polynomial,
confirming the finite nature of $h_1$ from the eventual periodicity of $g_1$.

More generally, $g_m(q)$ eventually becomes a quasi-polynomial in $q$ (i.e.\
its coefficients eventually satisfy a linear recurrence), and $(q)_m$ serves as
the ``minimal annihilator'' that converts $g_m$ into the polynomial $h_m$.

\subsection{Observations across profiles}

\begin{center}
\begin{tabular}{c|c|c|c|c}
$c$ & $d$ & $h_1$ & $g_1$ stabilises at & involution pairs for $n=1$ \\
\hline
$(2,1,1)$ & 4 & $3q + q^2 + q^3$ & $5$ (from $w \geq 3$) & 1 pair at $w=1$ \\
$(2,2,1)$ & 5 & $3q + 2q^2 + q^3 + q^4$ & $7$ (from $w \geq 4$) & 1 pair at $w=1$ \\
$(3,1,1)$ & 5 & $3q + q^2 + 2q^3 + q^5$ & $7$ (from $w \geq 5$) & 1 pair at $w=1$ \\
\end{tabular}
\end{center}

In all cases, $[q^1]\,h_1 = 3$ (corresponding to the 3 single-box CPPs) and the
involution pairs exactly one of these with the unique $j=1$ object.
The coefficient $[q^1]\,Q_1 = 3 - 1 = 2$ is universal across all nontrivial
profiles ($d \geq 2$) with $\gcd(d, 3) = 1$.


%% ========================================================================
\section{The cyclic descent injection}
\label{sec:cyclic-descent}
%% ========================================================================

In this section we describe a deterministic injection from weight-$(w{-}1)$
max-$1$ CPPs to weight-$w$ max-$1$ CPPs governed by the \emph{slack vector}.
When $d \not\equiv 0 \pmod{3}$, this injection is valid for all~$w$ and all
profiles~$c$ with $c_0 + c_1 + c_2 = d$, and its fixed points account for the
polynomial $h_1(q) = (q)_1 \cdot g_1(q)$.

\subsection{Slack vectors}

\begin{definition}[Slack vector]
For a max-$1$ CPP triple $(l_1, l_2, l_3)$ satisfying~\eqref{eq:length-conditions},
the \emph{slack vector} is
\begin{equation}
  s = (s_1, s_2, s_3) := (l_3 + c_0 - l_1,\; l_1 + c_1 - l_2,\; l_2 + c_2 - l_3).
  \label{eq:slack-def}
\end{equation}
Each $s_i \geq 0$ measures how far the $i$-th interlacing
condition~\eqref{eq:interlace1}--\eqref{eq:interlace3} is from being tight.
Note that $s_1 + s_2 + s_3 = c_0 + c_1 + c_2 = d$.
\end{definition}

The slack vector determines the ``local geometry'' of a triple on the
interlacing polytope: tight constraints ($s_i = 0$) correspond to facets,
and the interior point $s = (d/3, d/3, d/3)$ (when $3 \mid d$) is the unique
\emph{balanced slack}.

\begin{remark}[Effect of incrementing on slack]
\label{rem:slack-effect}
Incrementing $l_k$ by~$1$ changes the slack as follows:
\begin{align*}
  l_1 \to l_1 + 1\colon &\quad s_1 \to s_1 - 1,\; s_2 \to s_2 + 1,\; s_3 \to s_3, \\
  l_2 \to l_2 + 1\colon &\quad s_1 \to s_1,\; s_2 \to s_2 - 1,\; s_3 \to s_3 + 1, \\
  l_3 \to l_3 + 1\colon &\quad s_1 \to s_1 + 1,\; s_2 \to s_2,\; s_3 \to s_3 - 1.
\end{align*}
In particular, incrementing $l_k$ decreases $s_k$ by~$1$ and increases
$s_{k-1 \bmod 3}$ by~$1$ (with cyclic indexing $s_0 = s_3$).
The increment is valid (producing a new CPP triple) if and only if $s_k > 0$.
\end{remark}

\subsection{Cyclic descents}

\begin{definition}[Cyclic descent]
A \emph{cyclic descent} of a slack vector $s = (s_1, s_2, s_3)$ at position~$k$
(for $k \in \{1, 2, 3\}$) is a strict inequality
\[
  s_k > s_{k-1 \bmod 3},
\]
where the indices are taken cyclically: $s_0 := s_3$, $s_{-1} := s_2$.
We write $\operatorname{CycDes}(s) = \{k \in \{1,2,3\} : s_k > s_{k-1 \bmod 3}\}$.
\end{definition}

\begin{lemma}
\label{lem:cycdes-nonempty}
If $s = (s_1, s_2, s_3)$ with $s_1 + s_2 + s_3 = d$ and $s_i \geq 0$,
and the $s_i$ are not all equal, then $\operatorname{CycDes}(s) \neq \emptyset$.
Moreover, $|\operatorname{CycDes}(s)| \in \{1, 2\}$.
\end{lemma}

\begin{proof}
If $\operatorname{CycDes}(s) = \emptyset$, then $s_k \leq s_{k-1}$ cyclically:
$s_1 \leq s_3 \leq s_2 \leq s_1$, forcing $s_1 = s_2 = s_3$.
The second claim follows since the three differences
$s_1 - s_3$, $s_2 - s_1$, $s_3 - s_2$ sum to~$0$, so the number of
strictly positive ones equals the number of strictly negative ones,
hence $|\operatorname{CycDes}(s)|$ equals the number of strictly negative
cyclic differences, which is $1$ or~$2$.
\end{proof}

\subsection{The injection rule}

\begin{definition}[Cyclic descent injection]
\label{def:injection}
Let $3 \nmid d$. For a max-$1$ CPP triple $(l_1, l_2, l_3)$ with slack vector
$s = (s_1, s_2, s_3)$, define
\begin{equation}
  \phi(l_1, l_2, l_3) := (l_1, l_2, l_3) + e_k,
  \label{eq:phi-def}
\end{equation}
where $e_k$ is the $k$-th standard basis vector and
\[
  k := \begin{cases}
    \text{the unique element of } \operatorname{CycDes}(s)
      & \text{if } |\operatorname{CycDes}(s)| = 1, \\[4pt]
    \operatorname*{argmin}_{j \in \operatorname{CycDes}(s)} s_{j-1 \bmod 3}
      & \text{if } |\operatorname{CycDes}(s)| = 2,
  \end{cases}
\]
with ties in the second case broken by choosing the smaller~$k$.
\end{definition}

\noindent
Equivalently: among all positions $k$ with $s_k > s_{k-1}$, choose the one
whose ``predecessor slack'' $s_{k-1}$ is smallest (tightest constraint),
with ties broken by the smallest index.

\begin{theorem}
\label{thm:injection}
Let $3 \nmid d$. The map $\phi$ defined in~\eqref{eq:phi-def} is a valid
injection from the set of max-$1$ CPP triples at weight~$w$ to those at
weight~$w+1$, for every $w \geq 1$.
\end{theorem}

\begin{proof}[Computational verification]
We have verified this theorem by exhaustive computation for all
$\binom{d+2}{2}$ profiles $c$ with $c_0 + c_1 + c_2 = d$, for:
\begin{itemize}
  \item $d = 1$: 3 profiles, $w \leq 30$.
  \item $d = 2$: 6 profiles, $w \leq 30$.
  \item $d = 4$: 15 profiles, $w \leq 30$.
  \item $d = 5$: 21 profiles, $w \leq 30$.
  \item $d = 7$: 36 profiles, $w \leq 30$.
  \item $d = 8$: 45 profiles, $w \leq 20$.
\end{itemize}
In each case, the injection is valid (produces a valid triple, and is injective)
at every weight. The verification is implemented in \texttt{test\_cyclic\_rule.py}.
\end{proof}

\begin{remark}[Failure for $3 \mid d$]
When $3 \mid d$, the balanced slack $s = (d/3, d/3, d/3)$ has
$\operatorname{CycDes}(s) = \emptyset$, so the cyclic descent rule is undefined.
Moreover, no static slack-based rule can work in this case: the count of max-$1$
CPP triples oscillates with period~$3$ (e.g.\ $3, 3, 4, 3, 3, 4, \ldots$
for $d = 3$), so the transition from a ``spike'' weight to a ``trough'' weight
sends more triples than there are targets, making injection impossible.
A different mechanism (possibly involving triples of consecutive weight levels)
is needed for $3 \mid d$.
\end{remark}

\subsection{Fixed points and the polynomial $h_1$}

The \emph{fixed points} of the injection $\phi$ at weight~$w$ are the max-$1$
CPP triples at weight~$w$ that are not in the image of~$\phi$.

\begin{proposition}
\label{prop:fixed-pts}
Let $3 \nmid d$ and $w \geq 2$. The number of fixed points of $\phi$ at
weight~$w$ equals $[q^w]\,h_1(q)$.
At weight~$1$, the number of fixed points is $[q^1]\,h_1(q) + 1 = 4$
(since all 3 triples at $w = 1$ are unmatched, plus there are no predecessors
at $w = 0$).
\end{proposition}

\begin{proof}[Computational verification]
Verified against the independently computed $h_1$ for all profiles
$c$ with $d \in \{2, 4, 5, 7\}$, checking that the injection fixed-point
count at weight~$w$ matches $[q^w]\,h_1(q)$ for $w \geq 2$.
\end{proof}

The connection to the full involution for $Q_{n,c}(q)$ is as follows.
The simplest case is $n = 1$: here $Q_1 = h_1 - q \cdot h_0$ and the
injection $\phi$ on max-$1$ triples gives $h_1 = (1 - q) \cdot g_1$.
The injection accounts for the factor $(1-q)$, and the remaining subtraction
$-q \cdot h_0$ (the $j = 1$ term) cancels one of the three weight-$1$ fixed
points, yielding $[q^1]\,Q_1 = 3 - 1 = 2$.

\subsection{Example: $c = (2,1,1)$, $d = 4$}

The 15 slack vectors on the simplex $\{s_1 + s_2 + s_3 = 4,\; s_i \geq 0\}$
and their assigned components are:

\begin{center}
\small
\begin{tabular}{cccc|c|cc}
$s_1$ & $s_2$ & $s_3$ & $|\operatorname{CycDes}|$ &
$k = \phi$-comp. & from & to \\
\hline
4 & 0 & 0 & 1 & 1 & $(4,0,0) \to (3,0,1)$ & $s_1 \downarrow, s_3 \uparrow$ \\
3 & 1 & 0 & 1 & 1 & $(3,1,0) \to (2,1,1)$ & \\
3 & 0 & 1 & 2 & 3 & $(3,0,1) \to (3,1,0)$ & picks $s_2=0$ \\
2 & 2 & 0 & 1 & 1 & $(2,2,0) \to (1,2,1)$ & \\
2 & 1 & 1 & 1 & 1 & $(2,1,1) \to (1,1,2)$ & \\
2 & 0 & 2 & 2 & 3 & $(2,0,2) \to (2,1,1)$ & picks $s_2=0$ \\
1 & 3 & 0 & 2 & 1 & $(1,3,0) \to (0,3,1)$ & picks $s_3=0$ \\
1 & 2 & 1 & 1 & 2 & $(1,2,1) \to (2,1,1)$ & \\
1 & 1 & 2 & 1 & 3 & $(1,1,2) \to (1,2,1)$ & \\
1 & 0 & 3 & 1 & 3 & $(1,0,3) \to (1,1,2)$ & \\
0 & 4 & 0 & 1 & 2 & $(0,4,0) \to (1,3,0)$ & \\
0 & 3 & 1 & 1 & 2 & $(0,3,1) \to (1,2,1)$ & \\
0 & 2 & 2 & 1 & 2 & $(0,2,2) \to (1,1,2)$ & \\
0 & 1 & 3 & 2 & 2 & $(0,1,3) \to (1,0,3)$ & picks $s_1=0$ \\
0 & 0 & 4 & 1 & 3 & $(0,0,4) \to (0,1,3)$ & \\
\end{tabular}
\end{center}

The slack-to-slack transition graph has a single absorbing cycle of length~3:
\[
  (2,1,1) \xrightarrow{l_1} (1,1,2) \xrightarrow{l_3} (1,2,1) \xrightarrow{l_2} (2,1,1),
\]
corresponding to the cyclic rotation of the ``nearly balanced'' slacks.
All other slacks eventually reach this cycle, and the stable regime has
5 triples per weight with 0 fixed points.

At weights $1$--$3$, the counts are $3, 4, 5$ triples respectively.
The injection produces $1$ fixed point at $w = 2$ (since $4 - 3 = 1$)
and $1$ at $w = 3$ (since $5 - 4 = 1$), giving
$h_1 = 3q + q^2 + q^3$.

\subsection{Towards a proof}

A conceptual proof of Theorem~\ref{thm:injection} would need to show:
\begin{enumerate}
  \item \textbf{Validity:} For any slack $s$ with $s_1 + s_2 + s_3 = d$ and
    $3 \nmid d$, the cyclic descent $k$ satisfies $s_k > 0$, so incrementing
    $l_k$ does not violate any constraint. This follows from
    Lemma~\ref{lem:cycdes-nonempty} since $s_k > s_{k-1} \geq 0$
    implies $s_k \geq 1$.

  \item \textbf{Injectivity:} Different triples at weight~$w$ produce
    different targets at weight~$w+1$. This requires showing that the map
    $s \mapsto s'$ (where $s' = s + e_{k-1} - e_k$ in the cyclic notation)
    is injective on the set of slack vectors appearing at weight~$w$.
    Since the slack determines the triple uniquely (given the weight), and
    the slack transformation is determined by the rule, it suffices to show
    that the induced map on slacks is injective. This amounts to verifying
    that no two distinct slack vectors map to the same successor slack under
    the cyclic descent rule.
\end{enumerate}

The injectivity in point~(2) appears to be the key obstacle. Computer verification
confirms it for all tested values, but a uniform combinatorial argument remains to
be found. The cyclic symmetry $\Z/3\Z$ of the slack simplex
(rotating indices $1 \to 2 \to 3 \to 1$) plays a central role: the rule respects
this symmetry in the sense that it ``follows the cyclic descent'', and the
$3 \nmid d$ condition ensures the absence of a fixed point under the cyclic action.


%% ========================================================================
\section*{Acknowledgements}
%% ========================================================================

These notes were developed with computational assistance using Python for
multipartition enumeration and exact rational $q$-series arithmetic.

\begin{thebibliography}{99}
\bibitem{Warnaar2021} S.~O.~Warnaar, \emph{The $A_2$ Andrews--Gordon identities
  and cylindric partitions}, arXiv:2111.07550v3, 2023.
\bibitem{Langer2021} R.~Langer, \emph{Cylindric partitions, lambda determinants,
  and semicircular systems}, PhD thesis, University of Queensland, 2021.
\bibitem{GK1997} I.~M.~Gessel and C.~Krattenthaler, \emph{Cylindric partitions},
  Trans.\ Amer.\ Math.\ Soc., 349 (1997), 429--479.
\bibitem{Borodin2007} A.~Borodin, \emph{Periodic Schur process and cylindric
  partitions}, Duke Math.\ J., 140 (2007), 391--468.
\end{thebibliography}

\end{document}
